%!TEX root = ../main.tex

\section{Work Experience}
\resumeSubHeadingListStart

  \resumeSubheading
    {Meta Reality Labs}{September 2021 -- Present}
    {Research Scientist — Neural Interfaces}{New York, NY}
    \resumeItemListStart
      \resumeItem{Led data collection and model development for decoding continuous joint kinematics from wearable surface biosignals, enabling a real-time cursor-control BCI interface. Scaled motion-capture data collection rigs to collect training data from hundreds of human subjects, applying neural scaling laws to optimize data collection budgets and inform long-term modeling.}
      \resumeItem{Spearheaded research utilizing Variational Autoencoders (VAEs) for out-of-distribution (OOD) detection in biosignals data.}
      \resumeItem{Oversaw the zero-to-one engineering of distributed machine learning pipelines built on PyTorch and Hydra for large-scale experiment sweeps, benchmarking, and corpus generation.}
      \resumeItem{Conducted 50+ technical interviews to grow the team and directly managed and mentored over 10 machine learning interns and junior researchers.}
      \resumeItem{Led open-sourcing of the most diverse large-scale surface biosignal dataset to date and corresponding model training and evaluation pipelines (\href{https://github.com/facebookresearch/generic-neuromotor-interface}{GitHub}).}
    \resumeItemListEnd

  \resumeSubheading
    {University College London — Gatsby Computational Neuroscience Unit}{September 2016 -- May 2021}
    {PhD Researcher}{London, UK}
    \resumeItemListStart
      \resumeItem{Developed mathematical models to analyze circuit mechanisms of biological motor learning and credit assignment during closed-loop brain-computer interface (BCI) control.}
      \resumeItem{Built, simulated, and optimized recurrent neural networks (RNNs) to derive testable experimental predictions regarding low-dimensional neural manifold shifts during BCI control.}
      \resumeItem{Conducted event-related analyses of electrocorticographic (ECoG) recordings from rat motor cortex using behavioral kinematics extracted from video with computer vision techniques.}
    \resumeItemListEnd

\resumeSubHeadingListEnd